%% fig_surplus_analysis.tex
%% 余剰分析：UE と SO の比較，最適混雑課金による死荷重損失の解消
%% Usage: %% fig_surplus_analysis.tex
%% 余剰分析：UE と SO の比較，最適混雑課金による死荷重損失の解消
%% Usage: %% fig_surplus_analysis.tex
%% 余剰分析：UE と SO の比較，最適混雑課金による死荷重損失の解消
%% Usage: %% fig_surplus_analysis.tex
%% 余剰分析：UE と SO の比較，最適混雑課金による死荷重損失の解消
%% Usage: \input{assets/assignment/fig_surplus_analysis.tex}
%%
%% 模式図のパラメータ（すべて正規化された単位）
%%   需要曲線     D(x)   = 4 - 0.5x         （右下がり）
%%   私的限界費用  PMC(x) = 1 + 0.5x = t_a(x_a)
%%   社会的限界費用 SMC(x) = 1 + x   = d/dx[x * t_a(x_a)]
%%
%%   SO 点 (D = SMC): x_SO = 2,  t_SO = PMC(2) = 2,  価格(含む課金) = 3
%%   UE 点 (D = PMC): x_UE = 3,  t_UE = PMC(3) = 2.5
%%   最適課金 τ* = SMC(x_SO) - PMC(x_SO) = 3 - 2 = 1
%%   自由流旅行時間 t_0 = PMC(0) = SMC(0) = 1
%%   死荷重損失の三角形 頂点: SO点(2,3), UE点(3,2.5), SMC(x_UE)=(3,4)

\begin{tikzpicture}[
  every node/.style={font=\small},
  xscale=1.6,    %% 横方向の拡大率（x 軸 0-4）
  yscale=1.4     %% 縦方向の拡大率（y 軸 0-4.5）
]

%% =============================================================
%% 塗り潰し（背面レイヤー）
%% =============================================================

%% 死荷重損失「紫の三角形」:  SO点 → SMC(x_UE) → UE点 で囲まれる三角形
%% 境界: 上辺 = SMC 曲線 (2,3)→(3,4), 右辺 = 垂直線 (3,4)→(3,2.5),
%%       下辺 = D 曲線 (3,2.5)→(2,3)
\fill[violet!30] (2, 3) -- (3, 4) -- (3, 2.5) -- cycle;

%% =============================================================
%% 曲線（3 本）
%% =============================================================

%% 需要曲線  D(x) = 4 - 0.5x
\draw[very thick, blue!75!black]
  (0, 4) -- (3.6, 2.2)
  node[right] {$D$};

%% 私的限界費用  PMC(x) = t_a(x_a) = 1 + 0.5x
\draw[very thick, green!55!black]
  (0, 1) -- (3.6, 2.8)
  node[right] {$\text{PMC} = t_a(x_a)$};

%% 社会的限界費用  SMC(x) = 1 + x （= d/dx[x * t_a] = t_a + x_a t_a'）
\draw[very thick, red!65!black]
  (0, 1) -- (3.3, 4.3)
  node[right, align=left] {$\text{SMC} = t_a + x_a\, t_a'$};

%% =============================================================
%% 補助破線
%% =============================================================

%% x_SO の垂直破線
\draw[gray, dashed, thin] (2, 0) -- (2, 3);

%% x_UE の垂直破線
\draw[gray, dashed, thin] (3, 0) -- (3, 2.5);

%% t_SO = PMC(x_SO) = 2 の水平破線
\draw[gray, dashed, thin] (0, 2) -- (2, 2);

%% t_UE = PMC(x_UE) = 2.5 の水平破線
\draw[gray, dashed, thin] (0, 2.5) -- (3, 2.5);

%% =============================================================
%% 均衡点（塗り潰し丸）
%% =============================================================

%% SO 点（D と SMC の交点, x=2, y=3）
\fill (2, 3) circle[radius=2.3pt];

%% UE 点（D と PMC の交点, x=3, y=2.5）
\fill (3, 2.5) circle[radius=2.3pt];

%% PMC(x_SO) の点（課金の下端, x=2, y=2）
\fill[gray!60] (2, 2) circle[radius=1.8pt];

%% =============================================================
%% 軸（前面レイヤー）
%% =============================================================

\draw[thick, ->] (-0.15, 0) -- (4.0, 0) node[right] {$x_a$（交通量）};
\draw[thick, ->] (0, -0.15) -- (0, 4.7) node[above] {費用（旅行時間）};

%% 原点ラベル
\node[below left, font=\small] at (0, 0) {O};

%% =============================================================
%% 軸の目盛り・ラベル
%% =============================================================

%% x 軸
\draw[thin] (2, 0) -- (2, -0.1);
\draw[thin] (3, 0) -- (3, -0.1);
\node[below]          at (2, 0)   {$x_{SO}$};
\node[below]          at (3, 0)   {$x_{UE}$};

%% y 軸
\draw[thin] (-0.1, 1)   -- (0, 1);
\draw[thin] (-0.1, 2)   -- (0, 2);
\draw[thin] (-0.1, 2.5) -- (0, 2.5);

\node[left, font=\footnotesize] at (0, 1)   {$t_0$};
\node[left]                     at (0, 2)   {$t_{SO}$};
\node[left]                     at (0, 2.5) {$t_{UE}$};

%% =============================================================
%% τ* 注釈（最適課金額：PMC(x_SO)〜SMC(x_SO) の垂直両矢印）
%% =============================================================

\draw[<->, >=stealth, thick]
  (2.2, 2.0) -- (2.2, 3.0)
  node[midway, right, font=\footnotesize] {$\tau^*$};

%% =============================================================
%% 領域ラベル
%% =============================================================

%% 死荷重損失ラベル（三角形の内側）
\node[font=\footnotesize, align=center, violet!70!black]
  at (2.72, 3.15) {死荷重\\[-2pt]損失};

%% =============================================================
%% 均衡点のラベル
%% =============================================================

\node[above left, font=\footnotesize] at (2.2, 3.1)    {SO点};
\node[below right, font=\footnotesize] at (2.8, 2.4) {UE点};

\end{tikzpicture}

%%
%% 模式図のパラメータ（すべて正規化された単位）
%%   需要曲線     D(x)   = 4 - 0.5x         （右下がり）
%%   私的限界費用  PMC(x) = 1 + 0.5x = t_a(x_a)
%%   社会的限界費用 SMC(x) = 1 + x   = d/dx[x * t_a(x_a)]
%%
%%   SO 点 (D = SMC): x_SO = 2,  t_SO = PMC(2) = 2,  価格(含む課金) = 3
%%   UE 点 (D = PMC): x_UE = 3,  t_UE = PMC(3) = 2.5
%%   最適課金 τ* = SMC(x_SO) - PMC(x_SO) = 3 - 2 = 1
%%   自由流旅行時間 t_0 = PMC(0) = SMC(0) = 1
%%   死荷重損失の三角形 頂点: SO点(2,3), UE点(3,2.5), SMC(x_UE)=(3,4)

\begin{tikzpicture}[
  every node/.style={font=\small},
  xscale=1.6,    %% 横方向の拡大率（x 軸 0-4）
  yscale=1.4     %% 縦方向の拡大率（y 軸 0-4.5）
]

%% =============================================================
%% 塗り潰し（背面レイヤー）
%% =============================================================

%% 死荷重損失「紫の三角形」:  SO点 → SMC(x_UE) → UE点 で囲まれる三角形
%% 境界: 上辺 = SMC 曲線 (2,3)→(3,4), 右辺 = 垂直線 (3,4)→(3,2.5),
%%       下辺 = D 曲線 (3,2.5)→(2,3)
\fill[violet!30] (2, 3) -- (3, 4) -- (3, 2.5) -- cycle;

%% =============================================================
%% 曲線（3 本）
%% =============================================================

%% 需要曲線  D(x) = 4 - 0.5x
\draw[very thick, blue!75!black]
  (0, 4) -- (3.6, 2.2)
  node[right] {$D$};

%% 私的限界費用  PMC(x) = t_a(x_a) = 1 + 0.5x
\draw[very thick, green!55!black]
  (0, 1) -- (3.6, 2.8)
  node[right] {$\text{PMC} = t_a(x_a)$};

%% 社会的限界費用  SMC(x) = 1 + x （= d/dx[x * t_a] = t_a + x_a t_a'）
\draw[very thick, red!65!black]
  (0, 1) -- (3.3, 4.3)
  node[right, align=left] {$\text{SMC} = t_a + x_a\, t_a'$};

%% =============================================================
%% 補助破線
%% =============================================================

%% x_SO の垂直破線
\draw[gray, dashed, thin] (2, 0) -- (2, 3);

%% x_UE の垂直破線
\draw[gray, dashed, thin] (3, 0) -- (3, 2.5);

%% t_SO = PMC(x_SO) = 2 の水平破線
\draw[gray, dashed, thin] (0, 2) -- (2, 2);

%% t_UE = PMC(x_UE) = 2.5 の水平破線
\draw[gray, dashed, thin] (0, 2.5) -- (3, 2.5);

%% =============================================================
%% 均衡点（塗り潰し丸）
%% =============================================================

%% SO 点（D と SMC の交点, x=2, y=3）
\fill (2, 3) circle[radius=2.3pt];

%% UE 点（D と PMC の交点, x=3, y=2.5）
\fill (3, 2.5) circle[radius=2.3pt];

%% PMC(x_SO) の点（課金の下端, x=2, y=2）
\fill[gray!60] (2, 2) circle[radius=1.8pt];

%% =============================================================
%% 軸（前面レイヤー）
%% =============================================================

\draw[thick, ->] (-0.15, 0) -- (4.0, 0) node[right] {$x_a$（交通量）};
\draw[thick, ->] (0, -0.15) -- (0, 4.7) node[above] {費用（旅行時間）};

%% 原点ラベル
\node[below left, font=\small] at (0, 0) {O};

%% =============================================================
%% 軸の目盛り・ラベル
%% =============================================================

%% x 軸
\draw[thin] (2, 0) -- (2, -0.1);
\draw[thin] (3, 0) -- (3, -0.1);
\node[below]          at (2, 0)   {$x_{SO}$};
\node[below]          at (3, 0)   {$x_{UE}$};

%% y 軸
\draw[thin] (-0.1, 1)   -- (0, 1);
\draw[thin] (-0.1, 2)   -- (0, 2);
\draw[thin] (-0.1, 2.5) -- (0, 2.5);

\node[left, font=\footnotesize] at (0, 1)   {$t_0$};
\node[left]                     at (0, 2)   {$t_{SO}$};
\node[left]                     at (0, 2.5) {$t_{UE}$};

%% =============================================================
%% τ* 注釈（最適課金額：PMC(x_SO)〜SMC(x_SO) の垂直両矢印）
%% =============================================================

\draw[<->, >=stealth, thick]
  (2.2, 2.0) -- (2.2, 3.0)
  node[midway, right, font=\footnotesize] {$\tau^*$};

%% =============================================================
%% 領域ラベル
%% =============================================================

%% 死荷重損失ラベル（三角形の内側）
\node[font=\footnotesize, align=center, violet!70!black]
  at (2.72, 3.15) {死荷重\\[-2pt]損失};

%% =============================================================
%% 均衡点のラベル
%% =============================================================

\node[above left, font=\footnotesize] at (2.2, 3.1)    {SO点};
\node[below right, font=\footnotesize] at (2.8, 2.4) {UE点};

\end{tikzpicture}

%%
%% 模式図のパラメータ（すべて正規化された単位）
%%   需要曲線     D(x)   = 4 - 0.5x         （右下がり）
%%   私的限界費用  PMC(x) = 1 + 0.5x = t_a(x_a)
%%   社会的限界費用 SMC(x) = 1 + x   = d/dx[x * t_a(x_a)]
%%
%%   SO 点 (D = SMC): x_SO = 2,  t_SO = PMC(2) = 2,  価格(含む課金) = 3
%%   UE 点 (D = PMC): x_UE = 3,  t_UE = PMC(3) = 2.5
%%   最適課金 τ* = SMC(x_SO) - PMC(x_SO) = 3 - 2 = 1
%%   自由流旅行時間 t_0 = PMC(0) = SMC(0) = 1
%%   死荷重損失の三角形 頂点: SO点(2,3), UE点(3,2.5), SMC(x_UE)=(3,4)

\begin{tikzpicture}[
  every node/.style={font=\small},
  xscale=1.6,    %% 横方向の拡大率（x 軸 0-4）
  yscale=1.4     %% 縦方向の拡大率（y 軸 0-4.5）
]

%% =============================================================
%% 塗り潰し（背面レイヤー）
%% =============================================================

%% 死荷重損失「紫の三角形」:  SO点 → SMC(x_UE) → UE点 で囲まれる三角形
%% 境界: 上辺 = SMC 曲線 (2,3)→(3,4), 右辺 = 垂直線 (3,4)→(3,2.5),
%%       下辺 = D 曲線 (3,2.5)→(2,3)
\fill[violet!30] (2, 3) -- (3, 4) -- (3, 2.5) -- cycle;

%% =============================================================
%% 曲線（3 本）
%% =============================================================

%% 需要曲線  D(x) = 4 - 0.5x
\draw[very thick, blue!75!black]
  (0, 4) -- (3.6, 2.2)
  node[right] {$D$};

%% 私的限界費用  PMC(x) = t_a(x_a) = 1 + 0.5x
\draw[very thick, green!55!black]
  (0, 1) -- (3.6, 2.8)
  node[right] {$\text{PMC} = t_a(x_a)$};

%% 社会的限界費用  SMC(x) = 1 + x （= d/dx[x * t_a] = t_a + x_a t_a'）
\draw[very thick, red!65!black]
  (0, 1) -- (3.3, 4.3)
  node[right, align=left] {$\text{SMC} = t_a + x_a\, t_a'$};

%% =============================================================
%% 補助破線
%% =============================================================

%% x_SO の垂直破線
\draw[gray, dashed, thin] (2, 0) -- (2, 3);

%% x_UE の垂直破線
\draw[gray, dashed, thin] (3, 0) -- (3, 2.5);

%% t_SO = PMC(x_SO) = 2 の水平破線
\draw[gray, dashed, thin] (0, 2) -- (2, 2);

%% t_UE = PMC(x_UE) = 2.5 の水平破線
\draw[gray, dashed, thin] (0, 2.5) -- (3, 2.5);

%% =============================================================
%% 均衡点（塗り潰し丸）
%% =============================================================

%% SO 点（D と SMC の交点, x=2, y=3）
\fill (2, 3) circle[radius=2.3pt];

%% UE 点（D と PMC の交点, x=3, y=2.5）
\fill (3, 2.5) circle[radius=2.3pt];

%% PMC(x_SO) の点（課金の下端, x=2, y=2）
\fill[gray!60] (2, 2) circle[radius=1.8pt];

%% =============================================================
%% 軸（前面レイヤー）
%% =============================================================

\draw[thick, ->] (-0.15, 0) -- (4.0, 0) node[right] {$x_a$（交通量）};
\draw[thick, ->] (0, -0.15) -- (0, 4.7) node[above] {費用（旅行時間）};

%% 原点ラベル
\node[below left, font=\small] at (0, 0) {O};

%% =============================================================
%% 軸の目盛り・ラベル
%% =============================================================

%% x 軸
\draw[thin] (2, 0) -- (2, -0.1);
\draw[thin] (3, 0) -- (3, -0.1);
\node[below]          at (2, 0)   {$x_{SO}$};
\node[below]          at (3, 0)   {$x_{UE}$};

%% y 軸
\draw[thin] (-0.1, 1)   -- (0, 1);
\draw[thin] (-0.1, 2)   -- (0, 2);
\draw[thin] (-0.1, 2.5) -- (0, 2.5);

\node[left, font=\footnotesize] at (0, 1)   {$t_0$};
\node[left]                     at (0, 2)   {$t_{SO}$};
\node[left]                     at (0, 2.5) {$t_{UE}$};

%% =============================================================
%% τ* 注釈（最適課金額：PMC(x_SO)〜SMC(x_SO) の垂直両矢印）
%% =============================================================

\draw[<->, >=stealth, thick]
  (2.2, 2.0) -- (2.2, 3.0)
  node[midway, right, font=\footnotesize] {$\tau^*$};

%% =============================================================
%% 領域ラベル
%% =============================================================

%% 死荷重損失ラベル（三角形の内側）
\node[font=\footnotesize, align=center, violet!70!black]
  at (2.72, 3.15) {死荷重\\[-2pt]損失};

%% =============================================================
%% 均衡点のラベル
%% =============================================================

\node[above left, font=\footnotesize] at (2.2, 3.1)    {SO点};
\node[below right, font=\footnotesize] at (2.8, 2.4) {UE点};

\end{tikzpicture}

%%
%% 模式図のパラメータ（すべて正規化された単位）
%%   需要曲線     D(x)   = 4 - 0.5x         （右下がり）
%%   私的限界費用  PMC(x) = 1 + 0.5x = t_a(x_a)
%%   社会的限界費用 SMC(x) = 1 + x   = d/dx[x * t_a(x_a)]
%%
%%   SO 点 (D = SMC): x_SO = 2,  t_SO = PMC(2) = 2,  価格(含む課金) = 3
%%   UE 点 (D = PMC): x_UE = 3,  t_UE = PMC(3) = 2.5
%%   最適課金 τ* = SMC(x_SO) - PMC(x_SO) = 3 - 2 = 1
%%   自由流旅行時間 t_0 = PMC(0) = SMC(0) = 1
%%   死荷重損失の三角形 頂点: SO点(2,3), UE点(3,2.5), SMC(x_UE)=(3,4)

\begin{tikzpicture}[
  every node/.style={font=\small},
  xscale=1.6,    %% 横方向の拡大率（x 軸 0-4）
  yscale=1.4     %% 縦方向の拡大率（y 軸 0-4.5）
]

%% =============================================================
%% 塗り潰し（背面レイヤー）
%% =============================================================

%% 死荷重損失「紫の三角形」:  SO点 → SMC(x_UE) → UE点 で囲まれる三角形
%% 境界: 上辺 = SMC 曲線 (2,3)→(3,4), 右辺 = 垂直線 (3,4)→(3,2.5),
%%       下辺 = D 曲線 (3,2.5)→(2,3)
\fill[violet!30] (2, 3) -- (3, 4) -- (3, 2.5) -- cycle;

%% =============================================================
%% 曲線（3 本）
%% =============================================================

%% 需要曲線  D(x) = 4 - 0.5x
\draw[very thick, blue!75!black]
  (0, 4) -- (3.6, 2.2)
  node[right] {$D$};

%% 私的限界費用  PMC(x) = t_a(x_a) = 1 + 0.5x
\draw[very thick, green!55!black]
  (0, 1) -- (3.6, 2.8)
  node[right] {$\text{PMC} = t_a(x_a)$};

%% 社会的限界費用  SMC(x) = 1 + x （= d/dx[x * t_a] = t_a + x_a t_a'）
\draw[very thick, red!65!black]
  (0, 1) -- (3.3, 4.3)
  node[right, align=left] {$\text{SMC} = t_a + x_a\, t_a'$};

%% =============================================================
%% 補助破線
%% =============================================================

%% x_SO の垂直破線
\draw[gray, dashed, thin] (2, 0) -- (2, 3);

%% x_UE の垂直破線
\draw[gray, dashed, thin] (3, 0) -- (3, 2.5);

%% t_SO = PMC(x_SO) = 2 の水平破線
\draw[gray, dashed, thin] (0, 2) -- (2, 2);

%% t_UE = PMC(x_UE) = 2.5 の水平破線
\draw[gray, dashed, thin] (0, 2.5) -- (3, 2.5);

%% =============================================================
%% 均衡点（塗り潰し丸）
%% =============================================================

%% SO 点（D と SMC の交点, x=2, y=3）
\fill (2, 3) circle[radius=2.3pt];

%% UE 点（D と PMC の交点, x=3, y=2.5）
\fill (3, 2.5) circle[radius=2.3pt];

%% PMC(x_SO) の点（課金の下端, x=2, y=2）
\fill[gray!60] (2, 2) circle[radius=1.8pt];

%% =============================================================
%% 軸（前面レイヤー）
%% =============================================================

\draw[thick, ->] (-0.15, 0) -- (4.0, 0) node[right] {$x_a$（交通量）};
\draw[thick, ->] (0, -0.15) -- (0, 4.7) node[above] {費用（旅行時間）};

%% 原点ラベル
\node[below left, font=\small] at (0, 0) {O};

%% =============================================================
%% 軸の目盛り・ラベル
%% =============================================================

%% x 軸
\draw[thin] (2, 0) -- (2, -0.1);
\draw[thin] (3, 0) -- (3, -0.1);
\node[below]          at (2, 0)   {$x_{SO}$};
\node[below]          at (3, 0)   {$x_{UE}$};

%% y 軸
\draw[thin] (-0.1, 1)   -- (0, 1);
\draw[thin] (-0.1, 2)   -- (0, 2);
\draw[thin] (-0.1, 2.5) -- (0, 2.5);

\node[left, font=\footnotesize] at (0, 1)   {$t_0$};
\node[left]                     at (0, 2)   {$t_{SO}$};
\node[left]                     at (0, 2.5) {$t_{UE}$};

%% =============================================================
%% τ* 注釈（最適課金額：PMC(x_SO)〜SMC(x_SO) の垂直両矢印）
%% =============================================================

\draw[<->, >=stealth, thick]
  (2.2, 2.0) -- (2.2, 3.0)
  node[midway, right, font=\footnotesize] {$\tau^*$};

%% =============================================================
%% 領域ラベル
%% =============================================================

%% 死荷重損失ラベル（三角形の内側）
\node[font=\footnotesize, align=center, violet!70!black]
  at (2.72, 3.15) {死荷重\\[-2pt]損失};

%% =============================================================
%% 均衡点のラベル
%% =============================================================

\node[above left, font=\footnotesize] at (2.2, 3.1)    {SO点};
\node[below right, font=\footnotesize] at (2.8, 2.4) {UE点};

\end{tikzpicture}
